\begin{circuitikz}[
    american,
    voltage dir=RP,
    >=latex,
    %>={Latex[length=5mm, width=2mm]},
    alt={}
    ]

    % Define the number of rungs in the diagram
    \def\numRungs{5}

    % Calculate the length of the vertical power rails dynamically based on number of rungs
    \pgfmathsetmacro{\railLength}{-2 * \numRungs}

    %======================================================================================================
    % Component Grid
    % Spacing = 3x, 2y
    %======================================================================================================
    \foreach \i in {0,...,5} {
        \foreach \j in {0,...,20} { % Change to 0,...,19 if you need 20 individual Y points down to -40
        
        % Calculate spatial positions
        \pgfmathsetmacro{\px}{\i * 3}
        \pgfmathsetmacro{\py}{-1 - (\j * 2)}
        
        % Name coordinate using loop indices: (C-column-row) e.g., (C-0-0), (C-5-10)
        \coordinate (C-\i-\j) at (\px, \py);
        
        % OPTIONAL: Uncomment the 2 lines below to display grid labels while building
        % \fill[red] (C-\i-\j) circle (1.5pt);
        %\node[anchor=south west, scale=0.5, gray] at (C-\i-\j) {(\i,\j)};
        }
    }

    %======================================================================================================
    % Foreground Layer
    % Only contains the vertical power rails
    % Length of the rails is automatic, calculated based on the number of rungs defined above
    %======================================================================================================
    \begin{pgfonlayer}{ft}

        %======================================================================================================
        % Scope sets options for both lines - BakerBlack, very thick
        %======================================================================================================
        \begin{scope}[
            draw=BakerBlack,
            text=BakerBlack,
            color=BakerBlack,
            font=\small,
            ultra thick
        ]
        
            \draw (0,0) -- (0,\railLength);

            \draw (15,0) -- (15,\railLength);

        \end{scope}
    \end{pgfonlayer}

    %======================================================================================================
    % Main Layer
    % All actual logic goes here
    %======================================================================================================
    \begin{pgfonlayer}{main}
        
        %======================================================================================================
        % Scope sets options for both lines - BakerBlack for draw and text
        %======================================================================================================
        \begin{scope}[
            draw=BakerBlack,
            text=BakerBlack,
            color=BakerBlack,
            font=\small
        ]

            %=================================================================
            % Rung 1
            %=================================================================
            \draw (C-0-0) to[vC, l={\texttt{BlinkOn.DN}}] (C-1-0) -- (C-3-0) pic (TON1) {TON};
            \draw (TON1-OUT) -- (C-5-0);

            \node[anchor=south] at (TON1-TOP) {TON};
            \node[anchor=north] at (TON1-TOP) {\texttt{BlinkOff}};
            \node[anchor=west] at (TON1-PRE) {500};
            \node[anchor=west] at (TON1-ACC) {506};

            %=================================================================
            % Rung 3 (timer boxes take 2)
            %=================================================================
            \draw (C-0-2) to[C, l={\texttt{BlinkOff.DN}}] (C-1-2) to[vC, l={\texttt{BlinkOn.DN}}] (C-2-2) -- (C-3-2) pic (TON2) {TON};
            \draw (TON2-OUT) -- (C-5-2);

            \node[anchor=south] at (TON2-TOP) {TON};
            \node[anchor=north] at (TON2-TOP) {\texttt{BlinkOn}};
            \node[anchor=west] at (TON2-PRE) {500};
            \node[anchor=west] at (TON2-ACC) {503};

            \draw (C-0-4) to[C, l={\texttt{BlinkOff.DN}}] (C-1-4) -- (C-4-4) to[esource, l={\texttt{BlinkBit}}] (C-5-4);

        \end{scope}
    \end{pgfonlayer}

    % =================================================================
    % Background layer
    % =================================================================
    \begin{pgfonlayer}{bg}
        
        %==================================================================
        % Highlight Scope
        % This layer exists for highlighting true logic, the scope sets the
        % highlight options without having to repeat them every \draw
        % draw=UpNorthGreen, line width=10pt, opacity=0.4
        %==================================================================
        \begin{scope}[
            draw=UpNorthGreen,
            line width=10pt,
            opacity=0.4
        ]

            %==================================================================
            % Left vertical rung - leave this alone, it should ALWAYS be green
            % It automatically draws itself to the length defined at the top
            %==================================================================
            \draw (0,0) -- (0,\railLength);
        
            \draw (C-0-0) -- (C-3-0);

            \draw (TON1-EN1) -- (TON1-EN2);
            \draw (TON1-DN1) -- (TON1-DN2);

            \draw (TON2-EN1) -- (TON2-EN2);
            \draw (TON2-DN1) -- (TON2-DN2);

            \draw (C-0-2) -- (C-3-2);

            \draw (C-0-4) -- (C-5-4);

        \end{scope}
    \end{pgfonlayer}

\end{circuitikz}